\chapter{Lire un calendrier macroéconomique et un écran de taux}\label{ch:m2:reading-a-macro-calendar-and-a-rates-screen}

La plupart des vendredis en début de mois, à 08:30 heure de New York, le Bureau of Labor
Statistics publie le rapport sur l'emploi. De janvier 2023 à septembre 2026, le
rendement du Treasury à 2 ans a bougé, les 44 jours de publication, de 10.2 points de base en moyenne,
à la hausse comme à la baisse~; sur les 887 autres jours de négociation, de 4.6. Le marché des taux vit au rythme d'un calendrier de
telles publications et réunions de banques centrales, et un trader le lit comme un pilote lit une
carte météo~: où se trouvera la turbulence, à quel point elle peut devenir violente, et
à quel point sa position est exposée lorsqu'elle arrive. Ce chapitre referme le livre sur ce savoir-faire~: le calendrier,
le consensus et la surprise, le positionnement à l'approche des événements, les formes que peut prendre une courbe et
la façon de les négocier, et ce qu'un écran de taux montre en un coup d'œil.

\section{Le calendrier}

\begin{definition}[Publication économique]\label{def:m2:reading-a-macro-calendar-and-a-rates-screen:release}
Une \emph{publication économique}\index{publication économique} est la diffusion programmée d'une
statistique officielle, comme l'emploi ou les prix à la consommation, à une date et une heure annoncées
à l'avance par l'organisme qui la produit.
\end{definition}

Aux États-Unis, parmi les publications qui font bouger les taux, le rapport sur l'emploi et
l'indice des prix à la consommation sont tous deux publiés à 08:30, heure de l'Est~; les décisions
de politique monétaire de la Federal Reserve tombent à l'issue de ses huit réunions annuelles programmées. Le calendrier de
septembre 2026 montre comment elles s'articulent
(\cref{fig:m2:reading-a-macro-calendar-and-a-rates-screen:calendar})~: le rapport sur l'emploi d'août
publié le vendredi 4 septembre, l'indice des prix d'août le vendredi 11 septembre, et la
réunion du Federal Open Market Committee les 15 et 16 septembre, l'une des quatre réunions de l'année
assorties des projections économiques du Comité.

\begin{omfigure}
\centering
\begin{tikzpicture}[font=\small,x=0.62cm]
\fill[omThm!12] (5,-0.2) rectangle (18,0.2);
\draw[thick,-{Stealth[length=2.5mm]}] (0.5,0) -- (21.5,0);
\foreach \d in {1,2,3,4,5,6,7,8,9,10,11,12,13,14,15,16,17,18,19,20,21} \draw (\d,-0.08) -- (\d,0.08);
\foreach \d in {1,4,5,7,11,14,15,16,17,21} \node[font=\scriptsize,above] at (\d,0.2) {\d};
\node[font=\footnotesize,text=omThm] at (11.5,1.0) {Période de silence du FOMC~: du samedi 5 au jeudi 17 septembre};
\node[draw,thick,fill=omDef!10,rounded corners=2pt,align=center,font=\footnotesize] (p) at (4,-1.9) {Ven. 4, 08:30\\rapport sur l'emploi};
\node[draw,thick,fill=omProp!10,rounded corners=2pt,align=center,font=\footnotesize] (c) at (11,-1.9) {Ven. 11, 08:30\\IPC};
\node[draw,thick,fill=omThm!10,rounded corners=2pt,align=center,font=\footnotesize] (f) at (16.5,-1.9) {Mar. 15 -- Mer. 16\\FOMC, avec projections};
\draw[-{Stealth[length=2mm]}] (p.north) -- (4,-0.1);
\draw[-{Stealth[length=2mm]}] (c.north) -- (11,-0.1);
\draw[-{Stealth[length=2mm]}] (f.north) -- (16,-0.1);
\end{tikzpicture}

\omcaption{Le calendrier américain des taux pour les trois premières semaines de septembre 2026~: le
rapport sur l'emploi et l'IPC à 08:30, heure de l'Est, et la réunion du FOMC des 15 et 16
septembre et sa \omterm{def:m2:reading-a-macro-calendar-and-a-rates-screen:blackout}{période de silence}, du deuxième samedi précédant la réunion au
lendemain de celle-ci. Dates tirées des calendriers du BLS et de la Federal Reserve.}\label{fig:m2:reading-a-macro-calendar-and-a-rates-screen:calendar}
\end{omfigure}

\begin{definition}[Période de silence]\label{def:m2:reading-a-macro-calendar-and-a-rates-screen:blackout}
Une \emph{période de silence}\index{période de silence} est l'intervalle entourant une réunion de politique monétaire
au cours de laquelle les responsables d'une banque centrale s'abstiennent de tout commentaire public sur l'économie et
la politique monétaire.
\end{definition}

La règle du FOMC, réaffirmée en janvier 2026, fait débuter la \omterm{def:m2:reading-a-macro-calendar-and-a-rates-screen:blackout}{période de silence} à minuit le deuxième
samedi avant une réunion et la fait se terminer à 23:59 le lendemain~: pour une réunion du mardi--mercredi,
elle court du samedi dix jours plus tôt au jeudi. En septembre 2026, l'IPC
tombe dans cette période, si bien que le marché doit intégrer le rapport sans entendre les
membres du Comité. Les calendriers recensent aussi les jours fériés, les adjudications (chapitre~4) et les
rééquilibrages d'indices (chapitre~17)~; le calendrier d'un trader est l'union de ces éléments.

\begin{dated}{2026-09}{Calendrier américain des taux}\label{dat:m2:reading-a-macro-calendar-and-a-rates-screen:calendar}
Rapport sur l'emploi 2026~: 08:30 ET, par exemple 4 septembre, 2 octobre, 6 novembre, 4 décembre~; aucun
rapport n'a été publié en octobre 2025, pendant l'interruption des crédits budgétaires fédéraux. IPC 2026~:
08:30 ET, par exemple 11 septembre, 14 octobre. FOMC 2026~: 27--28 janvier, 17--18 mars, 28--29
avril, 16--17 juin, 28--29 juillet, 15--16 septembre, 27--28 octobre, 8--9 décembre (mars,
juin, septembre et décembre avec projections). Règle de \omterm{def:m2:reading-a-macro-calendar-and-a-rates-screen:blackout}{période de silence} du FOMC réaffirmée le 27 janvier 2026.
\end{dated}

\section{Consensus et surprise}

\begin{definition}[Prévision consensuelle, surprise statistique]\label{def:m2:reading-a-macro-calendar-and-a-rates-screen:surprise}
La \emph{prévision consensuelle}\index{prévision consensuelle} d'une publication est la médiane des
prévisions des économistes recueillies avant sa publication. La \emph{surprise statistique}\index{surprise statistique}
est la valeur publiée moins le consensus, souvent divisée par l'écart-type des
surprises passées pour rendre les publications comparables.
\end{definition}

Les marchés intègrent le consensus avant la publication~; ce qui les fait bouger, c'est la surprise. La
première mesure de l'importance d'une publication est l'ampleur des mouvements de marché les jours
où elle sort (\cref{fig:m2:reading-a-macro-calendar-and-a-rates-screen:moves}). Les jours de rapport sur l'emploi
de 2023 à 2026, le rendement à 2 ans a bougé 2.2 fois autant que les autres jours, celui à 10 ans 1.7
fois autant, et la pente 2s10s 1.4 fois autant~: le rapport porte sur la trajectoire de la politique monétaire, que
le segment court de la courbe l'intègre le plus directement.

\begin{omfigure}
\begin{tikzpicture}
\begin{axis}[width=11cm,height=4.8cm,ybar,bar width=18pt,
  ylabel={variation quotidienne absolue moyenne (pb)},
  tick label style={font=\small},label style={font=\small},
  xtick=data,xticklabels={2 ans,10 ans,2s10s},enlarge x limits=0.3,
  ymin=0,ymax=12,grid=major,grid style={gray!25},
  legend columns=2,legend style={font=\footnotesize,at={(0.5,-0.22)},anchor=north,draw=none,fill=none,/tikz/every even column/.append style={column sep=8pt}},
  nodes near coords,nodes near coords style={font=\scriptsize},
  every node near coord/.append style={/pgf/number format/fixed,/pgf/number format/fixed zerofill,/pgf/number format/precision=1},
  table/col sep=comma]
\addplot[fill=omDef!60,draw=omDef] table[x=k,y=payroll]{figdata/markets-2/31-reading-a-macro-calendar-and-a-rates-screen/eventmoves.csv};
\addplot[fill=gray!35,draw=gray] table[x=k,y=other]{figdata/markets-2/31-reading-a-macro-calendar-and-a-rates-screen/eventmoves.csv};
\legend{44 jours de rapport sur l'emploi,887 autres jours}
\end{axis}
\end{tikzpicture}

\omcaption{Variation quotidienne absolue moyenne des rendements du Treasury à 2 ans et à 10 ans, et de la
pente 2s10s, aux jours de rapport sur l'emploi et aux autres jours, de janvier 2023 à septembre 2026.
Données~: Federal Reserve H.15 via FRED (DGS2, DGS10)~; dates de publication tirées des
archives du BLS.}\label{fig:m2:reading-a-macro-calendar-and-a-rates-screen:moves}
\end{omfigure}

La seconde mesure est la sensibilité~: le nombre de points de base dont un rendement se déplace par unité de
\omterm{def:m2:reading-a-macro-calendar-and-a-rates-screen:surprise}{surprise statistique} standardisée, estimé en régressant les variations sur les surprises sur de nombreuses publications.
Les données de consensus sont vendues par des fournisseurs, si bien que le tutoriel présente la méthode sur des données synthétiques~:
avec 200 publications dont les sensibilités vraies sont de 8 et 5 points de base pour les rendements à 2 et à 10 ans,
la régression donne 8.4 et 5.1, avec des erreurs types de 0.2.

\section{Positionnement autour des événements}

Une position anodine un jour calme peut devenir dangereuse lors d'une \omterm{def:m2:reading-a-macro-calendar-and-a-rates-screen:release}{publication économique}. La question
avant un événement est l'ampleur du mouvement auquel on pourrait être exposé, à laquelle répond l'historique des jours d'événement,
et de savoir si l'on veut vraiment cette exposition~: un trader qui a une opinion l'exprime~; un
teneur de marché réduit son risque ou élargit ses cotations à l'approche du chiffre~; un portefeuille qui n'a pas vocation à
parier sur les données réduit sa duration ou couvre le court terme.

\begin{definition}[Steepener, flattener, butterfly de courbe]\label{def:m2:reading-a-macro-calendar-and-a-rates-screen:steep}
Un \emph{steepener}\index{steepener} est une position qui gagne lorsque la courbe des taux
se pentifie, typiquement longue d'une obligation plus courte et courte d'une plus longue, avec un DV01 égal~; un
\emph{flattener}\index{flattener} est la position inverse. Un \emph{butterfly de courbe}\index{butterfly de courbe} (curve fly) est
une position sur trois échéances, longue ou courte de l'échéance médiane contre les deux ailes, dimensionnée pour être
neutre au niveau et à la pente de la courbe~; son prix se mesure par deux fois le rendement médian
moins les rendements des deux ailes.
\end{definition}

Un \omterm{def:m2:reading-a-macro-calendar-and-a-rates-screen:steep}{steepener} neutre en DV01 est insensible aux mouvements parallèles~: il gagne ou perd la variation
de la pente multipliée par son DV01 par jambe. Avec le calculateur d'obligations du chapitre 3, un billet illustratif
à 10 ans, de coupon 4.875~\% et de rendement 4.96~\%, dont le DV01 vaut 0.0772 pour 100~; un billet
à 2 ans, de coupon 4.625~\% et de rendement 4.71~\%, dont le DV01 vaut 0.0183. Vendre 100~millions de dollars de l'échéance à 10 ans représente un
DV01 de 77\,168 dollars par point de base~; la jambe à 2 ans qui la compense correspond à 421~millions de dollars.

\section{Lire les formes de la courbe}

\begin{definition}[Bull steepening et bear flattening]\label{def:m2:reading-a-macro-calendar-and-a-rates-screen:bull}
Un mouvement de la courbe des taux est un \emph{bull steepening}\index{bull steepening} lorsque les rendements baissent
et que la courbe se pentifie, le court terme baissant davantage, et un \emph{bear flattening}\index{bear flattening}
lorsque les rendements montent et que la courbe s'aplatit, le court terme montant davantage. Le bull flattening et
le bear steepening sont les deux autres combinaisons.
\end{definition}

La plupart des jours de rapport sur l'emploi font bouger la courbe selon l'une des deux façons définies ci-dessus, le court terme
menant le mouvement (\cref{fig:m2:reading-a-macro-calendar-and-a-rates-screen:scatter})~: sur les 44 jours,
14 ont été des \omterm{def:m2:reading-a-macro-calendar-and-a-rates-screen:bull}{bull steepening} et 20 des \omterm{def:m2:reading-a-macro-calendar-and-a-rates-screen:bull}{bear flattening}. Le 2 août 2024, le rendement à 2 ans a reculé de 28 points de base
et celui à 10 ans de 19~; le 4 octobre 2024, le rendement à 2 ans a monté de 23 et celui à 10 ans de 13. Un rapport faible
tire vers le bas les taux directeurs anticipés, ce qui fait baisser le court terme le plus fortement~; un rapport fort
produit l'effet inverse. Les mouvements des deux autres types, où le long terme mène, renvoient souvent
à autre chose que la trajectoire de la politique monétaire~: prime de terme, offre, anticipations d'inflation.

\begin{omfigure}
\begin{tikzpicture}
\begin{axis}[width=9cm,height=7.2cm,
  xlabel={variation du rendement à 2 ans (pb)},ylabel={variation du rendement à 10 ans (pb)},
  tick label style={font=\small},label style={font=\small},
  xmin=-32,xmax=32,ymin=-26,ymax=26,grid=major,grid style={gray!25},
  legend columns=2,legend style={font=\footnotesize,at={(0.5,-0.2)},anchor=north,draw=none,fill=none,/tikz/every even column/.append style={column sep=8pt}},
  scatter/classes={bull-steepening={mark=*,omDef},bear-flattening={mark=square*,omThm},bull-flattening={mark=triangle*,omProp},bear-steepening={mark=diamond*,gray}},
  table/col sep=comma]
\addplot[scatter,only marks,mark size=2pt,scatter src=explicit symbolic] table[x=d2,y=d10,meta=kind]{figdata/markets-2/31-reading-a-macro-calendar-and-a-rates-screen/payrollmoves.csv};
\addplot[black,thin,dashed,domain=-26:26,forget plot] {x};
\addplot[gray,thin,forget plot] coordinates {(-32,0) (32,0)};
\addplot[gray,thin,forget plot] coordinates {(0,-26) (0,26)};
\legend{bull steepening,bear flattening,bull flattening,bear steepening}
\end{axis}
\end{tikzpicture}

\omcaption{Variations des rendements du Treasury à 2 et à 10 ans lors des 44 jours de rapport sur l'emploi de
janvier 2023 à septembre 2026, classés selon la moyenne des deux variations (bull ou
bear) et la variation de 2s10s (pentification au-dessus de la diagonale en pointillés, aplatissement en dessous).
Le \omterm{def:m2:reading-a-macro-calendar-and-a-rates-screen:bull}{bull steepening}, à gauche au-dessus de la diagonale, et le \omterm{def:m2:reading-a-macro-calendar-and-a-rates-screen:bull}{bear flattening}, à droite en dessous,
sont au nombre de 34 sur les 44. Données~: H.15 via FRED~; dates de publication du BLS.}\label{fig:m2:reading-a-macro-calendar-and-a-rates-screen:scatter}
\end{omfigure}

\section{Un écran de taux d'un coup d'œil}

L'écran d'un trader en taux condense le marché en quelques chiffres~: le niveau de la courbe~;
ses pentes (2s10s, 5s30s)~; sa courbure, le \omterm{def:m2:the-fx-options-market:rr}{butterfly} 2s5s10s~; la \omterm{def:m2:short-term-interest-rate-futures:path}{trajectoire implicite des taux directeurs} du
chapitre 8~; les \omterm{def:m2:interest-rate-swaps:spread}{spreads de swap} et la \omterm{def:m2:fx-swaps-forwards-and-the-cross-currency-basis:basis}{base de devises} des chapitres 9 et 16~; la volatilité
du chapitre 13~; et le calendrier. Pris ensemble, ils indiquent où se situe le marché, ce qu'il
anticipe et ce qui a changé. De septembre 2025 à septembre 2026, la courbe du Treasury
a monté et s'est aplatie~: le rendement à 2 ans est passé de 3.61~\% à 4.71~\%, celui à 10 ans de 4.15~\% à 4.96~\%,
et 2s10s a baissé de 54 à 25 points de base, tandis que le \omterm{def:m2:the-fx-options-market:rr}{butterfly} 2s5s10s est passé de $-34$ à $-1$
(\cref{fig:m2:reading-a-macro-calendar-and-a-rates-screen:curves}).

\begin{omfigure}
\begin{tikzpicture}
\begin{axis}[width=11cm,height=5cm,
  xlabel={échéance (années)},ylabel={rendement (\%)},
  tick label style={font=\small},label style={font=\small},
  xmode=log,log basis x=10,xtick={2,5,10,30},xticklabels={2,5,10,30},
  xmin=1.7,xmax=35,ymin=3.4,ymax=5.5,grid=major,grid style={gray!25},
  legend columns=3,legend style={font=\footnotesize,at={(0.5,-0.3)},anchor=north,draw=none,fill=none,/tikz/every even column/.append style={column sep=8pt}},
  table/col sep=comma]
\addplot[gray,very thick,mark=*] table[x=tenor,y=d20250922]{figdata/markets-2/31-reading-a-macro-calendar-and-a-rates-screen/curves.csv};
\addplot[omProp,very thick,mark=*] table[x=tenor,y=d20260323]{figdata/markets-2/31-reading-a-macro-calendar-and-a-rates-screen/curves.csv};
\addplot[omDef,very thick,mark=*] table[x=tenor,y=d20260922]{figdata/markets-2/31-reading-a-macro-calendar-and-a-rates-screen/curves.csv};
\legend{22 septembre 2025,23 mars 2026,22 septembre 2026}
\end{axis}
\end{tikzpicture}

\omcaption{La courbe du Treasury à 2, 5, 10 et 30 ans à trois dates espacées de six mois~:
sur l'année, elle a monté et s'est aplatie, un \omterm{def:m2:reading-a-macro-calendar-and-a-rates-screen:bull}{bear flattening} mené par le rendement à 2 ans. Données~: Federal
Reserve H.15 via FRED (DGS2, DGS5, DGS10, DGS30).}\label{fig:m2:reading-a-macro-calendar-and-a-rates-screen:curves}
\end{omfigure}

\begin{dated}{2026-09}{Un écran du Treasury}\label{dat:m2:reading-a-macro-calendar-and-a-rates-screen:screen}
Rendements à échéance constante le 22 septembre 2026~: 2 ans 4.71~\%, 5 ans 4.83~\%, 10 ans 4.96~\%,
30 ans 5.29~\%~; 2s10s 25 points de base, 5s30s 46, \omterm{def:m2:the-fx-options-market:rr}{butterfly} 2s5s10s $-1$. Un an plus tôt~: 3.61, 3.71,
4.15 et 4.77~\%~; 54, 106 et $-34$.
\end{dated}

\section{Tutoriel~: un calendrier de publications et un écran de courbe}\label{tut:m2:reading-a-macro-calendar-and-a-rates-screen:macrocal}

\begin{tutorial}
\textbf{Objectif.} Construire un calendrier de publications avec la \omterm{def:m2:reading-a-macro-calendar-and-a-rates-screen:blackout}{période de silence}, mesurer les mouvements les jours d'événement,
estimer les sensibilités aux surprises, synthétiser la courbe et dimensionner une transaction neutre en DV01.
\textbf{État final~:} \cref{fig:m2:reading-a-macro-calendar-and-a-rates-screen:calendar,fig:m2:reading-a-macro-calendar-and-a-rates-screen:moves,fig:m2:reading-a-macro-calendar-and-a-rates-screen:scatter,fig:m2:reading-a-macro-calendar-and-a-rates-screen:curves}
et les chiffres du problème du week-end.
\begin{tutsteps}
\item \textbf{Le calendrier}~: événements, règle de la \omterm{def:m2:reading-a-macro-calendar-and-a-rates-screen:blackout}{période de silence} et indicateurs.
\omcode{code/firm/macrocal/firm_macrocal.py}{13}{29}{Événements, période de silence du FOMC et calendrier.}\label{lst:m2:reading-a-macro-calendar-and-a-rates-screen:cal}
\item \textbf{L'écran}~: pentes et \omterm{def:m2:reading-a-macro-calendar-and-a-rates-screen:steep}{butterfly de courbe}, classification des mouvements, mouvements les jours d'événement,
dimensionnement neutre en DV01 et P\&L.
\omcode{code/firm/macrocal/firm_macrocal.py}{32}{61}{Synthèse de la courbe, classification, mouvements les jours d'événement et dimensionnement.}\label{lst:m2:reading-a-macro-calendar-and-a-rates-screen:screen}
\item \textbf{Surprises}~: la sensibilité des rendements aux \omterm{def:m2:reading-a-macro-calendar-and-a-rates-screen:surprise}{surprises statistiques} standardisées.
\omcode{code/firm/macrocal/firm_macrocal.py}{64}{72}{Sensibilité aux surprises par moindres carrés.}\label{lst:m2:reading-a-macro-calendar-and-a-rates-screen:sens}
\item \textbf{Exécuter} \texttt{macro\_demo.event\_stats()}, \texttt{macro\_demo.problem()},
\texttt{macro\_demo.surprise\_regression()} et \texttt{fig\_macro.py}.
\end{tutsteps}
\textbf{Pour aller plus loin.} Ajouter les dates de publication de l'IPC et comparer les jours d'IPC aux jours de créations d'emplois~;
puis estimer la sensibilité du 2 ans aux surprises sur les créations d'emplois si l'on dispose de données de consensus,
et la comparer au mouvement moyen.
\end{tutorial}

\section{Construction~: calendrier des publications et synthèse de la courbe}\label{bld:m2:reading-a-macro-calendar-and-a-rates-screen:macrocal}

\begin{build}
\bfield{Rôle} La note du matin de la petite société et ses limites de risque autour des événements~: quelles publications sortent à quel moment,
quelles réunions sont en \omterm{def:m2:reading-a-macro-calendar-and-a-rates-screen:blackout}{période de silence}, à quoi ressemble la courbe et quelle est la taille des
positions sur courbe, exprimée en DV01.
\bfield{Interface} \texttt{Event(day, time, name)}~; \texttt{blackout(meeting\_start,
meeting\_end)}~; \texttt{calendar(events, meetings)}~; \texttt{curve\_summary(y2, y5, y10,
y30)}~; \texttt{classify(d2, d10)}~; \texttt{event\_day\_moves(dates, y, events)}~;
\texttt{dv01\_neutral(dv01\_short\_leg, dv01\_long\_leg, notional\_short)}~;
\texttt{steepener\_pnl(dv01, d2\_bp, d10\_bp)}~; \texttt{sensitivity(surprises, moves)}.
\bfield{Règles} Rendements en pourcentage, variations en points de base~; \omterm{def:m2:reading-a-macro-calendar-and-a-rates-screen:blackout}{période de silence} du deuxième samedi
précédant la réunion jusqu'au lendemain de la réunion (exception des jours fériés fédéraux non implémentée)~;
P\&L au premier ordre pour les opérations sur courbe~; MCO avec ordonnée à l'origine pour les sensibilités.
\bfield{Tests d'acceptation} \texttt{code/firm/macrocal/tests/}~: les \omterm{def:m2:reading-a-macro-calendar-and-a-rates-screen:blackout}{périodes de silence} de septembre 2026 et de janvier
2026~; le calendrier signale l'IPC dans la \omterm{def:m2:reading-a-macro-calendar-and-a-rates-screen:blackout}{période de silence} et les créations d'emplois en dehors~; pentes et
\omterm{def:m2:reading-a-macro-calendar-and-a-rates-screen:steep}{butterfly de courbe} du 22 septembre 2026~; les quatre classifications~; moyennes des jours d'événement, notionnel neutre en DV01
et P\&L du \omterm{def:m2:reading-a-macro-calendar-and-a-rates-screen:steep}{steepener} sur de petits cas~; la régression retrouve une pente connue.
\bfield{Pour aller plus loin} Séries de consensus et de surprises issues d'un fournisseur de données~; mouvements intrajournaliers autour des
publications~; la \omterm{def:m2:short-term-interest-rate-futures:path}{trajectoire implicite des taux directeurs} avant et après chaque réunion (chapitre 8)~; un écran du matin
qui combine les constructions de tous les chapitres.
\end{build}

\begin{omsources}
\item Bureau of Labor Statistics, calendriers de publication de l'Employment Situation
et de l'IPC~; archive des publications de l'Employment Situation.
\item Federal Reserve, calendriers du FOMC~; FOMC Policy on External Communications of Committee
Participants (réaffirmée en janvier 2026). \item Board of Governors of the Federal Reserve System, H.15 selected interest rates, via
FRED (DGS2, DGS5, DGS10, DGS30).
\end{omsources}

\section{Exercices}

\begin{exercise}[$\star$]\label{exo:m2:reading-a-macro-calendar-and-a-rates-screen:1}
Le FOMC se réunit le mardi 27 et le mercredi 28 octobre 2026. Quand commence sa \omterm{def:m2:reading-a-macro-calendar-and-a-rates-screen:blackout}{période de silence}
et quand se termine-t-elle~?
\end{exercise}

\begin{exercise}[$\star$]\label{exo:m2:reading-a-macro-calendar-and-a-rates-screen:2}
Le rendement du 2 ans baisse de 10 points de base et celui du 10 ans de 4. Nommer le mouvement.
\end{exercise}

\begin{exercise}[$\star$]\label{exo:m2:reading-a-macro-calendar-and-a-rates-screen:3}
Calculer le 2s10s, le 5s30s et le \omterm{def:m2:the-fx-options-market:rr}{butterfly} 2s5s10s le 22 septembre 2025.
\end{exercise}

\begin{exercise}[$\star\star$]\label{exo:m2:reading-a-macro-calendar-and-a-rates-screen:4}
Pourquoi le rapport sur l'emploi fait-il davantage bouger le rendement du 2 ans que celui du 10 ans~?
\end{exercise}

\begin{exercise}[$\star\star$]\label{exo:m2:reading-a-macro-calendar-and-a-rates-screen:5}
Un \omterm{def:m2:reading-a-macro-calendar-and-a-rates-screen:steep}{steepener} neutre en DV01 a un DV01 de USD~50\,000 par jambe. Que rapporte-t-il si le 2 ans
baisse de 15 points de base et le 10 ans de 5~? Si les deux baissent de 10~?
\end{exercise}

\begin{exercise}[$\star\star$]\label{exo:m2:reading-a-macro-calendar-and-a-rates-screen:6}
Pourquoi un pass-through lors d'une publication, sans position, est-il parfois la meilleure position~?
\end{exercise}

\begin{exercise}[$\star\star\star$]\label{exo:m2:reading-a-macro-calendar-and-a-rates-screen:7}
\emph{Programmation.} Avec \texttt{sensitivity} et les surprises synthétiques, combien de publications
faut-il pour que l'erreur type de la sensibilité du 2 ans tombe à 0.5 point de base~?
Vérifier la réponse par simulation.
\end{exercise}

\begin{exercise}[$\star\star\star$]\label{exo:m2:reading-a-macro-calendar-and-a-rates-screen:8}
\emph{Repérer l'erreur.} «~Le 2 ans bouge de 10 points de base les jours de créations d'emplois~; ma position a une
DV01 de USD~100\,000~; mon risque sur les créations d'emplois est donc de USD~1 million.~» La corriger.
\end{exercise}

\section{Problème~: créations d'emplois le vendredi}

\begin{problem}[{Problème du week-end --- un steepener avant le rapport sur l'emploi}]\label{pb:m2:reading-a-macro-calendar-and-a-rates-screen:1}
Un jeudi soir, un trader anticipe un rapport sur l'emploi décevant et veut une position sur courbe
qui en profite sans parier sur le niveau des taux. Il vend USD~100 millions d'un
\omterm{def:m2:money-markets:bill}{bon du Trésor} illustratif à 10 ans (coupon 4.875~\%, rendement 4.96~\%) et achète un bon à 2 ans (coupon
4.625~\%, rendement 4.71~\%), neutre en DV01, avec règlement le 22 septembre 2026.

\medskip\noindent\textbf{Partie I --- Dimensionnement.}
\begin{enumerate}
\item Quelles sont les DV01 des deux bons pour 100 de valeur nominale~?
\item Quel est le DV01 de la jambe à 10 ans~?
\item Quelle part du 2 ans le trader achète-t-il~?
\item Pourquoi neutre en DV01 plutôt qu'à notionnel égal~?
\item Que gagne ou perd la position si tous les rendements bougent de la même quantité~?
\end{enumerate}

\medskip\noindent\textbf{Partie II --- Le rapport.}
\begin{enumerate}[resume]
\item Supposons que la journée se répète le 2 août 2024~: le 2 ans baisse de 28 points de base et le 10 ans
de 19. Quel est le P\&L au premier ordre~? Et après repricing complet~?
\item Supposons plutôt le 4 octobre 2024~: le 2 ans monte de 23 et le 10 ans de 13.
\item Nommer chacun des deux mouvements.
\item À quelle fréquence, sur les 44 jours de créations d'emplois de 2023--2026, la courbe a-t-elle bougé selon l'un de ces deux
sens~?
\item D'où vient l'écart entre le P\&L au premier ordre et le P\&L après repricing~?
\end{enumerate}

\medskip\noindent\textbf{Partie III --- Le calendrier.}
\begin{enumerate}[resume]
\item Le rapport sort le 4 septembre 2026. Le FOMC est-il en \omterm{def:m2:reading-a-macro-calendar-and-a-rates-screen:blackout}{période de silence}~?
\item L'indice des prix à la consommation sort une semaine plus tard. Pourquoi la \omterm{def:m2:reading-a-macro-calendar-and-a-rates-screen:blackout}{période de silence} compte-t-elle pour la manière dont le marché l'accueille~?
\item De combien le 2 ans bouge-t-il davantage les jours de créations d'emplois que les autres jours~?
\item Qu'y a-t-il d'autre au calendrier qui pourrait faire bouger la position cette semaine-là~?
\item Comment dimensionneriez-vous la position pour perdre au plus 500\,000~USD lors d'une
journée type de mauvaises créations d'emplois~?
\end{enumerate}

\medskip\noindent\textbf{Partie IV --- Jugement.}
\begin{enumerate}[resume]
\item Pourquoi le \omterm{def:m2:reading-a-macro-calendar-and-a-rates-screen:steep}{steepener} pourrait-il perdre de l'argent même si le rapport est faible~?
\item Que surveilleriez-vous à l'écran avant et après la publication~?
\item Comment la transaction se rapporte-t-elle à la \omterm{def:m2:short-term-interest-rate-futures:path}{trajectoire implicite des taux directeurs} du chapitre~8~?
\item Énoncer le \emph{résultat nommé}~: la position \omterm{def:m2:reading-a-macro-calendar-and-a-rates-screen:steep}{steepener} 2s10s neutre en DV01 et son P\&L sous
chacun des deux mouvements.
\item En une phrase~: sur quoi parie un \omterm{def:m2:reading-a-macro-calendar-and-a-rates-screen:steep}{steepener} neutre en DV01~?
\end{enumerate}
\end{problem}

\section{Questions d'entretien}

\begin{interviewq}[$\star$ \iqroles{trader}]\label{iq:m2:reading-a-macro-calendar-and-a-rates-screen:1}
Quelles sont les publications américaines les plus importantes pour les taux, et à quel moment sortent-elles~?
\end{interviewq}

\begin{interviewq}[$\star$ \iqroles{trader, chercheur}]\label{iq:m2:reading-a-macro-calendar-and-a-rates-screen:2}
Définir le \omterm{def:m2:reading-a-macro-calendar-and-a-rates-screen:bull}{bull steepening} et le \omterm{def:m2:reading-a-macro-calendar-and-a-rates-screen:bull}{bear flattening}, et donner une cause de chacun.
\end{interviewq}

\begin{interviewq}[$\star\star$ \iqroles{trader}]\label{iq:m2:reading-a-macro-calendar-and-a-rates-screen:3}
Comment construit-on un \omterm{def:m2:reading-a-macro-calendar-and-a-rates-screen:steep}{steepener} 2s10s, et quels sont ses risques~?
\end{interviewq}

\begin{interviewq}[$\star\star$ \iqroles{chercheur}]\label{iq:m2:reading-a-macro-calendar-and-a-rates-screen:4}
Comment mesureriez-vous la sensibilité du marché à une \omterm{def:m2:reading-a-macro-calendar-and-a-rates-screen:release}{publication économique}~?
\end{interviewq}

\begin{interviewq}[$\star\star$ \iqroles{trader, risque}]\label{iq:m2:reading-a-macro-calendar-and-a-rates-screen:5}
Que mettriez-vous sur un écran de taux d'une page, et pourquoi~?
\end{interviewq}

\begin{interviewq}[$\star\star\star$ \iqroles{développeur}]\label{iq:m2:reading-a-macro-calendar-and-a-rates-screen:6}
Concevoir le système qui ingère une publication à 08:30:00 et met à jour, pour chaque desk, le risque et
les cotations en moins d'une seconde.
\end{interviewq}
